\chapter{Fine-Needle Aspiration Biopsy of the Thyroid}

\section*{Overview}
Thyroid fine-needle aspiration (FNA) biopsy uses a thin needle to collect cells from a thyroid nodule for cytologic examination.

\section*{Why This Test or Procedure Is Ordered}
It helps determine whether a thyroid nodule is benign, malignant, or indeterminate and guides decisions about surveillance, molecular testing, or surgery.

\section*{How This Test or Procedure Is Performed}
The neck is cleaned and ultrasound is used to guide one or more needle passes into the nodule. Local anesthetic may be used. The samples are sent to a pathology laboratory.

\section*{Risks and Results}
Temporary discomfort, bruising, bleeding, and rarely infection or injury to nearby structures can occur. Results may be benign, malignant, nondiagnostic, or indeterminate and must be interpreted with ultrasound and clinical findings.

\section*{References}
\begin{enumerate}
  \item MedlinePlus. Thyroid Fine-Needle Aspiration Biopsy. U.S. National Library of Medicine; 2025.
  \item Current Medical Diagnosis and Treatment. McGraw-Hill; 2025.
\end{enumerate}
\clearpage
